\chapter{Returns and Volatility}
\label{ch:returns}

\begin{companioncode}{qf-04-returns}
Code: \texttt{crates/qf-04-returns/}\\
Dependencies: \crate{qf-common}, \crate{qf-03-stocks}\\
Run: \texttt{cargo run -p qf-04-returns --example returns\_analysis}
\end{companioncode}

\section{Learning Objectives}

This chapter introduces the foundational vocabulary of quantitative finance.
After completing it you can:

\begin{itemize}
    \item Compute simple and logarithmic returns and know when to use each
    \item Measure volatility and annualise it for cross-frequency comparison
    \item Evaluate risk-adjusted performance with the Sharpe and Sortino ratios
    \item Quantify drawdowns and recovery time from a price series
    \item Use the \texttt{Returns} trait to abstract return computation across
      data sources
\end{itemize}

\begin{keyinsight}
Returns, not prices, are the currency of quant research. Prices are
non-stationary; returns are (approximately) stationary. Every quant model in
this book starts by transforming prices into returns.
\end{keyinsight}

\section{Simple vs Log Returns}

\subsection{Definitions}

\marginnote{%
\textbf{Simple return}\\[0.3em]
$r_t = \dfrac{P_t - P_{t-1}}{P_{t-1}}$\\[0.5em]
\textbf{Log return}\\[0.3em]
$r_t = \ln\!\left(\dfrac{P_t}{P_{t-1}}\right)$
}

The simple (arithmetic) return is the percentage change in price:
\[
r_t = \frac{P_t - P_{t-1}}{P_{t-1}}
\]
The log (continuously compounded) return is the natural log of the price ratio:
\[
r_t = \ln\!\left(\frac{P_t}{P_{t-1}}\right) = \ln(P_t) - \ln(P_{t-1})
\]

The two are related exactly by
\[
r^{\text{log}}_t = \ln\!\left(1 + r^{\text{simple}}_t\right),
\]
and for small returns ($|r| \ll 1$) they are nearly equal.

\subsection{When to use which}

\begin{center}
\begin{tabular}{lcc}
\toprule
\textbf{Property} & \textbf{Simple} & \textbf{Log} \\
\midrule
Interpretation & ``I made 10\%'' & Less intuitive \\
Additivity over time & No & Yes \\
Additivity over assets & Yes & No \\
Symmetry & Symmetric & Unbounded below \\
Compounding & Multiplicative & Additive \\
\bottomrule
\end{tabular}
\end{center}

\begin{keyinsight}
Log returns are additive across time: the $k$-period log return is the sum of
the $k$ single-period log returns. This makes them natural for time-series
modelling. Simple returns are additive across assets in a portfolio (weighted
by portfolio weights), which makes them natural for portfolio returns.
\end{keyinsight}

\subsection{Rust implementation}

\begin{lstlisting}[language=Rust]
/// Simple return: (P_t - P_{t-1}) / P_{t-1}
pub fn simple_returns(prices: &[f64]) -> Vec<f64> {
    if prices.len() < 2 { return Vec::new(); }
    let mut out = Vec::with_capacity(prices.len() - 1);
    for i in 1..prices.len() {
        let prev = prices[i - 1];
        out.push(if prev == 0.0 { 0.0 }
                 else { (prices[i] - prev) / prev });
    }
    out
}

/// Log return: ln(P_t / P_{t-1})
pub fn log_returns(prices: &[f64]) -> Vec<f64> {
    if prices.len() < 2 { return Vec::new(); }
    let mut out = Vec::with_capacity(prices.len() - 1);
    for i in 1..prices.len() {
        let (prev, curr) = (prices[i - 1], prices[i]);
        out.push(if prev > 0.0 && curr > 0.0 { (curr / prev).ln() }
                 else { 0.0 });
    }
    out
}
\end{lstlisting}

\marginnote{%
\textbf{Cumulative return}\\[0.3em]
$R_{0\to t} = \prod_{k=1}^{t}(1+r_k) - 1$
}

The cumulative return over a period is the running compounding of
single-period simple returns:
\[
R_{0 \to t} = \prod_{k=1}^{t}(1 + r_k) - 1.
\]

\begin{lstlisting}[language=Rust]
pub fn cumulative_returns(returns: &[f64]) -> Vec<f64> {
    let mut out = Vec::with_capacity(returns.len());
    let mut wealth = 1.0_f64;
    for &r in returns {
        wealth *= 1.0 + r;
        out.push(wealth - 1.0);
    }
    out
}
\end{lstlisting}

\section{Volatility}

\marginnote{%
\textbf{Sample std dev}\\[0.3em]
$\sigma = \sqrt{\dfrac{1}{n-1}\sum_{t=1}^{n}(r_t - \bar r)^2}$
}

Volatility is the standard deviation of returns. We use the \emph{sample}
estimator (denominator $n-1$), consistent with \func{qf\_common::compute\_stats}:
\[
\sigma = \sqrt{\frac{1}{n-1}\sum_{t=1}^{n}(r_t - \bar r)^2}.
\]

\subsection{Annualisation}

\marginnote{%
\textbf{Annualised vol}\\[0.3em]
$\sigma_{\text{ann}} = \sigma_{\text{period}} \cdot \sqrt{m}$\\[0.3em]
where $m$ is the number of periods per year (252 for daily).
}

To compare volatilities across frequencies (daily, weekly, monthly), we
annualise. Under the assumption that returns are i.i.d., variance scales
linearly with the horizon, so standard deviation scales with the square root:
\[
\sigma_{\text{annual}} = \sigma_{\text{period}} \cdot \sqrt{m}.
\]

\begin{lstlisting}[language=Rust]
pub fn volatility(returns: &[f64]) -> f64 {
    let n = returns.len();
    if n < 2 { return 0.0; }
    let mean: f64 = returns.iter().sum::<f64>() / n as f64;
    let ssd: f64 = returns.iter()
        .map(|&r| { let d = r - mean; d * d })
        .sum();
    (ssd / (n - 1) as f64).sqrt()
}

pub fn annualized_volatility(returns: &[f64], periods_per_year: f64) -> f64 {
    if periods_per_year <= 0.0 { return 0.0; }
    volatility(returns) * periods_per_year.sqrt()
}
\end{lstlisting}

\begin{warning}
The $\sqrt{m}$ rule assumes returns are independent and identically
distributed. For assets with autocorrelated or volatility-clustered returns
(see Chapter~\ref{ch:volatility}), annualisation by $\sqrt{m}$ overstates the
true annual volatility.
\end{warning}

\subsection{Rolling volatility}

A rolling window reveals how volatility changes over time---useful for
regime detection and risk monitoring:

\begin{lstlisting}[language=Rust]
pub fn rolling_volatility(returns: &[f64], window: usize) -> Vec<f64> {
    if window == 0 || window > returns.len() { return Vec::new(); }
    let mut out = Vec::with_capacity(returns.len() - window + 1);
    for i in window - 1..returns.len() {
        let slice = &returns[i + 1 - window..=i];
        out.push(volatility(slice));
    }
    out
}
\end{lstlisting}

\section{Risk-Adjusted Performance}

\subsection{Sharpe ratio}

\marginnote{%
\textbf{Sharpe ratio}\\[0.3em]
$S = \dfrac{\E[r] - r_f}{\sigma(r)}$
}

The Sharpe ratio measures excess return per unit of total risk:
\[
S = \frac{\bar r - r_f}{\sigma(r)}.
\]
To annualise, multiply by $\sqrt{m}$ (because both the mean excess return and
the standard deviation scale together under i.i.d.):
\[
S_{\text{annual}} = S_{\text{period}} \cdot \sqrt{m}.
\]

\begin{lstlisting}[language=Rust]
pub fn sharpe_ratio(returns: &[f64], risk_free_rate: f64) -> f64 {
    let n = returns.len();
    if n < 2 { return 0.0; }
    let vol = volatility(returns);
    if vol == 0.0 { return 0.0; }
    let mean: f64 = returns.iter().sum::<f64>() / n as f64;
    (mean - risk_free_rate) / vol
}

pub fn annualized_sharpe(returns: &[f64], rf: f64, periods_per_year: f64) -> f64 {
    if periods_per_year <= 0.0 { return 0.0; }
    sharpe_ratio(returns, rf) * periods_per_year.sqrt()
}
\end{lstlisting}

\subsection{Sortino ratio}

\marginnote{%
\textbf{Downside deviation}\\[0.3em]
$\sigma_D = \sqrt{\dfrac{1}{n}\sum_{r_t < r_f}(r_f - r_t)^2}$
}

The Sortino ratio replaces total volatility with \emph{downside} volatility
only. Upside volatility is good---it should not penalise the ratio:
\[
\text{Sortino} = \frac{\bar r - r_f}{\sigma_D}, \qquad
\sigma_D = \sqrt{\frac{1}{n}\sum_{r_t < r_f}(r_f - r_t)^2}.
\]

\begin{lstlisting}[language=Rust]
pub fn sortino_ratio(returns: &[f64], risk_free_rate: f64) -> f64 {
    let n = returns.len();
    if n < 2 { return 0.0; }
    let downside_sq: f64 = returns.iter()
        .map(|&r| {
            let short = risk_free_rate - r;
            if short > 0.0 { short * short } else { 0.0 }
        })
        .sum();
    if downside_sq == 0.0 { return 0.0; }
    let downside_dev = (downside_sq / n as f64).sqrt();
    let mean: f64 = returns.iter().sum::<f64>() / n as f64;
    (mean - risk_free_rate) / downside_dev
}
\end{lstlisting}

\begin{keyinsight}
Sortino $>$ Sharpe whenever the return distribution is right-skewed (more
upside than downside). For symmetric distributions the two ratios are equal in
expectation. Sortino is preferred for evaluating hedge-fund-style strategies
where upside volatility is desirable.
\end{keyinsight}

\section{Drawdown}

A drawdown is the percentage decline from a running peak:
\[
D_t = \frac{P_t - \max_{k \le t} P_k}{\max_{k \le t} P_k} \le 0.
\]

\marginnote{%
\textbf{Max drawdown}\\[0.3em]
$\text{MaxDD} = \min_t D_t$\\[0.3em]
The worst peak-to-trough decline over the period.
}

\begin{lstlisting}[language=Rust]
pub fn drawdown(prices: &[f64]) -> Vec<f64> {
    let mut out = Vec::with_capacity(prices.len());
    let mut running_max = f64::NEG_INFINITY;
    for &p in prices {
        if p > running_max { running_max = p; }
        out.push(if running_max == 0.0 { 0.0 }
                 else { (p - running_max) / running_max });
    }
    out
}

pub fn max_drawdown(prices: &[f64]) -> f64 {
    drawdown(prices).iter().copied().fold(0.0, f64::min)
}
\end{lstlisting}

The \texttt{DrawdownStats} struct summarises the worst episode:
\begin{itemize}
    \item \texttt{max\_drawdown}: the most negative drawdown value
    \item \texttt{max\_drawdown\_duration}: longest run of consecutive
      negative-drawdown observations
    \item \texttt{recovery\_time}: steps from the trough back to a new high
      (\texttt{None} if the series never recovered)
\end{itemize}

\begin{keyinsight}
Max drawdown measures \emph{path risk}, not just terminal risk. A strategy
that returns 10\% with a 50\% drawdown along the way is very different from one
that returns 10\% smoothly. Most investors abandon strategies with deep
drawdowns long before recovery---a risk that terminal-return statistics
completely miss.
\end{keyinsight}

\section{The Returns Trait}

\marginnote{%
\textbf{Trait benefits}\\[0.3em]
Define the capability (\emph{``can produce returns''}) once; implement for
slices, owned vectors, and OHLCV data.
}

Following the workspace-wide trait-first architecture rule, we abstract return
computation over price sources:

\begin{lstlisting}[language=Rust]
pub trait Returns {
    fn simple_returns(&self) -> Vec<f64>;
    fn log_returns(&self) -> Vec<f64>;
}

impl Returns for &[f64] { /* delegate to free functions */ }
impl Returns for Vec<f64> { /* delegate to the slice impl */ }
impl Returns for Vec<Ohlcv> {
    fn simple_returns(&self) -> Vec<f64> {
        let closes: Vec<f64> = self.iter().map(|o| o.close).collect();
        simple_returns(&closes)
    }
    // log_returns similarly uses closing prices
}
\end{lstlisting}

This lets generic code consume returns from any source that implements
\texttt{Returns}, keeping analytics decoupled from data loading.

\section{Example Output}

\begin{lstlisting}[language=bash]
$ cargo run -p qf-04-returns --example returns_analysis
Returns Analysis
================

Simple Returns: mean=0.0031, std=0.0215
Log Returns: mean=0.0029, std=0.0214
Volatility (daily): 0.0215
Volatility (annualized): 0.3417
Sharpe Ratio (annualized): 2.29
Sortino Ratio: 0.22
Max Drawdown: -7.50%
Drawdown duration: 6 steps, recovery: Some(6)

Trait check (Vec<f64>): first simple return = 0.0030
\end{lstlisting}

\section{Exercises}

\begin{enumerate}
    \item \textbf{Return conversion}: Write a function that converts a series
      of simple returns to log returns and vice versa. Verify the identity
      $r^{\text{log}} = \ln(1 + r^{\text{simple}})$ with a property test.

    \item \textbf{Annualisation sanity check}: For daily, weekly, and monthly
      returns of the same asset, verify that the annualised volatilities are
      approximately equal (they differ only by sampling noise).

    \item \textbf{Calmar ratio}: Implement the Calmar ratio: annualised return
      divided by the absolute value of the maximum drawdown. Compare it with
      the Sharpe ratio on a synthetic strategy.

    \item \textbf{Rolling Sharpe}: Implement a rolling Sharpe ratio (windowed
      mean and windowed std) and visualise how it changes over time. Where
      does it dip? Why?

    \item \textbf{Drawdown with recovery}: Extend \texttt{drawdown\_stats} to
      report the average drawdown (mean of negative drawdown values) and the
      proportion of time spent in drawdown.
\end{enumerate}

\section{Key Takeaways}

\begin{itemize}
    \item \textbf{Log returns} are time-additive; \textbf{simple returns} are
      portfolio-additive. Choose accordingly.
    \item \textbf{Volatility} uses sample std ($n-1$); annualise by multiplying
      by $\sqrt{m}$, valid under i.i.d.
    \item \textbf{Sharpe} penalises all volatility; \textbf{Sortino} penalises
      only downside.
    \item \textbf{Max drawdown} measures path risk---the deepest hole on the
      equity curve.
    \item \textbf{The \texttt{Returns} trait} keeps analytics decoupled from
      data sources, following the workspace trait-first rule.
\end{itemize}

\vspace{1em}
\noindent\textbf{Next chapter}: Your first backtest---an SMA crossover
strategy with P\&L accounting, transaction costs, and drawdown-based risk
evaluation.